% ------------------------- ESTADO DEL ARTE -------------------------
% Contenido (sumario de la planificación, §6): Agentes LLM, evaluación de agentes, simulación de usuarios, LLM-as-a-Judge,
% taxonomías de fallos, herramientas existentes y brecha identificada.
\clearpage
\blank
\tituloapa[Estado del arte]{ESTADO DEL ARTE}
\begingroup\dobleesp
\vspace{6pt}
Este capítulo presenta los conceptos en los que se apoya el trabajo, las herramientas y los trabajos académicos más cercanos a la propuesta, y la brecha que la motiva.
\par

\subtituloapa{Conceptos principales}
Un agente basado en modelos de lenguaje (LLM) es un sistema en el que un modelo interactúa con herramientas externas para completar tareas de varios pasos. La literatura de diseño distingue dos patrones: los \textit{workflows}, cuyo flujo de control está predefinido en código, y los agentes propiamente dichos, en los que el modelo dirige su propio proceso en un bucle con retroalimentación del entorno \cite{anthropic2024agents}. La distinción importa aquí en dos sentidos: el sistema que se evalúa es un agente, y el propio AgentProof se diseña como un workflow que delega la cognición en puntos acotados, como se detalla en el capítulo de la propuesta.
\par\vspace{6pt}
La evaluación de agentes se apoya en dos técnicas centrales \cite{yehudai2025survey}. La primera es la simulación de usuarios: un modelo de lenguaje interpreta al usuario y sostiene conversaciones de prueba con el agente, lo que permite ejercitar interacciones de varios turnos sin que una persona intervenga en cada ejecución. La segunda es el patrón \textit{LLM-as-a-Judge}, en el que un modelo califica ejecuciones contra criterios definidos. La confiabilidad de este evaluador automático no puede darse por supuesta y debe validarse contra jueces humanos \cite{yehudai2025survey}; la del usuario simulado se discute más adelante, junto con los trabajos que la han medido.
\par\vspace{6pt}
Para caracterizar cómo fallan los agentes se dispone de taxonomías específicas. La del AI Red Team de Microsoft cubre agentes individuales que usan herramientas \cite{msredteam2025taxonomy}. Trabajos recientes documentan además fallos silenciosos en sistemas en producción: errores cuya señal nunca llega a una persona de forma accionable y que, en su variante más peligrosa, el propio modelo convierte en una respuesta fluida y verosímil para el usuario \cite{wu2026narratives}. Estas taxonomías dan la semántica de las degradaciones controladas que se emplean en la evaluación experimental de este trabajo.
\par
\endgroup
